% abstract.tex -- Abstract
Training-free protein sequence generation via stochastic attention produces plausible family
members from small alignments, but treats all stored patterns equally and cannot direct generation
toward a user-specified functional subset---for example, sequences known from experimental screens
to bind a particular target. We show that assigning multiplicity weights to stored patterns
defines a tilted Hopfield energy whose Langevin dynamics sample from a known distribution,
requiring only a log-multiplicity bias on the softmax logits and no retraining. A single scalar
multiplicity ratio continuously interpolates between unconditioned generation and hard subset
curation. We characterize the phase transition as a function of the effective pattern count and
identify a family-dependent calibration gap between energy-level conditioning and decoded sequence
phenotype. This gap is predicted quantitatively by the geometric separation of functional subsets
in the principal component space of the alignment: well-separated families achieve high
phenotype transfer via multiplicity weighting alone, while overlapping families require hard
curation. Experiments on five Pfam families---Kunitz, SH3, WW, Homeobox, and Forkhead
domains---validate the theory across a fourfold range of separation indices, and extending
the analysis to omega-conotoxin peptides confirms the monotonic trend. Applied to
omega-conotoxin peptides targeting a voltage-gated calcium channel implicated in pain
signaling, curated memory seeding from a small set of characterized binders produces over a
thousand novel candidates that preserve the primary pharmacophore and all experimentally
identified binding determinants. These results establish stochastic attention as a practical
small-set amplifier for hit expansion from functional screens, enabling practitioners to convert
a handful of experimentally characterized sequences into diverse candidate libraries without
retraining a generative model.
